% =====================================================================
%  Workshop 4 -- ORM with FastAPI, SQLAlchemy and MySQL
%  Web Applications -- Semester II 2026
%
%  COMO USAR ESTE ARCHIVO (leeme):
%
%  1. Este .tex compila TAL CUAL, sin ninguna imagen. Donde va cada
%     captura aparece una caja gris con el nombre del archivo que
%     espera y la descripcion de lo que debe mostrar.
%
%  2. Para insertar una captura NO hay que editar nada: basta con
%     guardar la imagen dentro de la carpeta  img/  con el nombre
%     exacto que indica la caja gris (por ejemplo  img/fig01-show-tables.png ).
%     Al recompilar, la caja se reemplaza sola por la imagen.
%     Se aceptan las extensiones .png, .jpg, .jpeg y .pdf.
%
%  3. Los datos, salidas y SQL que aparecen en el documento son los
%     reales de la ejecucion del taller, no inventados.
%
%  4. Compilar con pdfLaTeX (Overleaf: New Project -> Upload Project,
%     subiendo este .tex y la carpeta img/).
% =====================================================================

\documentclass[10pt,a4paper]{article}

\usepackage[utf8]{inputenc}
\usepackage[T1]{fontenc}
\usepackage[english]{babel}
\usepackage[margin=2.1cm]{geometry}
\usepackage{graphicx}
\usepackage{xcolor}
\usepackage{listings}
\usepackage{booktabs}
\usepackage{float}
\usepackage{caption}
\usepackage{enumitem}
\usepackage[hidelinks]{hyperref}

\setlist{nosep,leftmargin=1.1cm}
\setlength{\parskip}{0.35em}

% ---------------------------------------------------------------------
%  Colores y estilos de codigo
% ---------------------------------------------------------------------
\definecolor{codebg}{RGB}{248,248,250}
\definecolor{codeframe}{RGB}{212,212,218}
\definecolor{codekey}{RGB}{0,84,166}
\definecolor{codestr}{RGB}{163,21,21}
\definecolor{codecom}{RGB}{106,115,125}
\definecolor{phfill}{RGB}{246,246,248}
\definecolor{phline}{RGB}{170,170,178}

\lstdefinestyle{base}{
  backgroundcolor=\color{codebg},
  frame=single,
  rulecolor=\color{codeframe},
  basicstyle=\ttfamily\scriptsize,
  keywordstyle=\color{codekey}\bfseries,
  stringstyle=\color{codestr},
  commentstyle=\color{codecom}\itshape,
  showstringspaces=false,
  breaklines=true,
  columns=fullflexible,
  keepspaces=true,
  xleftmargin=5pt,
  xrightmargin=5pt,
  aboveskip=6pt,
  belowskip=6pt,
}

\lstdefinestyle{py}{style=base,language=Python}
\lstdefinestyle{sql}{style=base,language=SQL,morekeywords={DISTINCT,INNER}}
\lstdefinestyle{sh}{style=base,language=bash}
\lstdefinestyle{plain}{style=base,language=}
\lstset{style=plain}

% ---------------------------------------------------------------------
%  Marcador de figura: si la imagen existe la inserta, si no dibuja
%  una caja vacia con la descripcion de la captura que falta.
% ---------------------------------------------------------------------
\newcommand{\phbox}[4]{%
  \begingroup
  \setlength{\fboxsep}{0pt}%
  \fcolorbox{phline}{phfill}{%
    \begin{minipage}[c][#2][c]{\dimexpr#1-2\fboxrule\relax}
      \centering
      {\ttfamily\scriptsize\color{phline}[ img/#3 ]}\\[0.7em]
      \begin{minipage}{0.88\linewidth}
        \centering\scriptsize\itshape #4
      \end{minipage}
    \end{minipage}%
  }%
  \endgroup
}

% #1 ancho  #2 alto de la caja  #3 nombre base del archivo  #4 descripcion
\newcommand{\shot}[4]{%
  \IfFileExists{img/#3.png}{\includegraphics[width=#1]{img/#3.png}}{%
  \IfFileExists{img/#3.jpg}{\includegraphics[width=#1]{img/#3.jpg}}{%
  \IfFileExists{img/#3.jpeg}{\includegraphics[width=#1]{img/#3.jpeg}}{%
  \IfFileExists{img/#3.pdf}{\includegraphics[width=#1]{img/#3.pdf}}{%
  \phbox{#1}{#2}{#3}{#4}}}}}%
}

\captionsetup{font=small,labelfont=bf,skip=4pt}

% ---------------------------------------------------------------------

\title{\vspace{-1.4cm}\textbf{Workshop 4}\\[0.2em]
\large ORM with FastAPI, SQLAlchemy and MySQL}
\author{Ismael Paredes\\[0.1em]
\normalsize Web Applications --- Semester II 2026 --- Prof. Francisco Hidrobo}
\date{September 19, 2026}

\begin{document}

\maketitle
\vspace{-0.8cm}

\section{General Information}

\begin{table}[H]
\centering
\footnotesize
\begin{tabular}{@{}ll@{\hspace{2cm}}ll@{}}
\toprule
Student & Ismael Paredes & Python & 3.12.10 \\
Course & Web Applications --- Sem. II 2026 & MySQL Server & 8.0.46 \\
Workshop & No. 4 --- ORM with FastAPI & SQLAlchemy & 2.0.54 (PyMySQL 1.2.3) \\
Date & September 19, 2026 & FastAPI & 0.141.1 (Uvicorn 0.53.0) \\
OS & Windows 11 Home (10.0.26200) & Database & \texttt{orm\_workshop} / \texttt{ormuser} \\
\bottomrule
\end{tabular}
\end{table}
\vspace{-0.3cm}
\noindent\footnotesize Repository: \url{https://github.com/CipherTuring/Web-Applications}\normalsize

\section{Objective}

The goal is to build a REST API that stores and retrieves information from a
relational database \emph{without writing SQL by hand}.

An \textbf{Object-Relational Mapper (ORM)} maps two worlds that do not
naturally fit together: the object world of Python (classes, instances,
attributes, references between objects) and the relational world of MySQL
(tables, rows, columns, foreign keys). SQLAlchemy defines that correspondence
once, in the model classes, and from then on the program manipulates ordinary
Python objects: creating an instance, assigning an attribute, appending to a
list or handing an object to \texttt{db.delete()}. The \texttt{INSERT},
\texttt{UPDATE} and \texttt{DELETE} statements that reach MySQL are generated
by the ORM. The schema itself is never written as SQL either: the four tables
were created directly from the Python classes.

\noindent
\begin{minipage}[t]{0.42\textwidth}
\begin{lstlisting}[style=plain]
Python Client / curl
        |
        | HTTP + JSON
        v
     FastAPI
        |
        | SQLAlchemy ORM
        v
      MySQL
\end{lstlisting}
\end{minipage}\hfill
\begin{minipage}[t]{0.55\textwidth}
\vspace{0.3cm}
\textbf{FastAPI} exposes the HTTP endpoints and validates the JSON crossing the
boundary. \textbf{SQLAlchemy} maps classes to tables, tracks the changes made to
the loaded objects and emits the SQL. \textbf{MySQL} stores the data and
enforces the foreign keys. The Python client is a separate program that knows
nothing about the ORM or the database: only URLs and JSON.

The application models an order system with four related entities:
\texttt{Customer}, \texttt{Order}, \texttt{OrderItem} and \texttt{Product}.
\end{minipage}

\section{Database and ORM Model}

\noindent
\begin{minipage}[t]{0.40\textwidth}
\begin{lstlisting}[style=plain]
Customer
   |
   +----< Order
             |
             +----< OrderItem >---- Product
\end{lstlisting}
\end{minipage}\hfill
\begin{minipage}[t]{0.57\textwidth}
\vspace{0.2cm}
\footnotesize
Every entity uses an auto-increment \texttt{id} as primary key. The foreign
keys are \texttt{orders.customer\_id}, \texttt{order\_items.order\_id} and
\texttt{order\_items.product\_id}. \texttt{OrderItem} is the association
between an order and a product.
\normalsize
\end{minipage}

\begin{table}[H]
\centering\footnotesize
\begin{tabular}{@{}p{2.1cm}p{5.4cm}p{8.2cm}@{}}
\toprule
\textbf{Entity} & \textbf{Purpose and attributes} & \textbf{Relationships} \\
\midrule
\texttt{Customer} & Places orders. \texttt{name}, \texttt{email}
(\texttt{unique=True}) &
One-to-many with \texttt{Order} (\texttt{orders}), \textbf{cascade}
\texttt{all, delete-orphan} \\
\addlinespace
\texttt{Order} & A purchase. \texttt{order\_date} (default today),
\texttt{status} (default \texttt{pending}) &
Belongs to one \texttt{Customer}; one-to-many with \texttt{OrderItem}
(\texttt{items}), \textbf{cascade} \texttt{all, delete-orphan} \\
\addlinespace
\texttt{OrderItem} & One line of an order. \texttt{quantity},
\texttt{unit\_price} &
Belongs to one \texttt{Order}, references one \texttt{Product} \\
\addlinespace
\texttt{Product} & Catalogue item. \texttt{name}, \texttt{price},
\texttt{stock} &
One-to-many with \texttt{OrderItem}, \textbf{no cascade} \\
\bottomrule
\end{tabular}
\end{table}

\texttt{unit\_price} is stored on the item deliberately: it is a snapshot of the
price at the moment of the purchase, so a later price change does not rewrite
the history of past orders.

\begin{table}[H]
\centering\footnotesize
\begin{tabular}{@{}p{3.2cm}p{12.6cm}@{}}
\toprule
\textbf{Construct} & \textbf{Purpose in this implementation} \\
\midrule
\texttt{ForeignKey} & Declares the link at the \emph{database} level; it is what
MySQL enforces and what appears as a \texttt{MUL} key. It works with values:
\texttt{orders.customer\_id} holds an integer. \\
\texttt{relationship()} & Declares the link at the \emph{object} level:
\texttt{order.customer} returns a \texttt{Customer}, \texttt{order.items} a list
of \texttt{OrderItem}. It turns a join into attribute access. \\
\texttt{back\_populates} & Pairs two \texttt{relationship()} declarations as the
two sides of one link and keeps them in sync: appending to \texttt{order.items}
also sets \texttt{item.order}. \\
\texttt{cascade="all,} \texttt{delete-orphan"} & Propagates session operations
from parent to children (\texttt{all}) and deletes a child removed from the
parent's collection (\texttt{delete-orphan}). \\
\bottomrule
\end{tabular}
\end{table}

Exactly two relationships use cascade: \texttt{Customer.orders} and
\texttt{Order.items}. \texttt{Product.order\_items} deliberately does not. The
criterion is ownership: an order \emph{owns} its items and a customer
\emph{owns} their orders, whereas an item only \emph{references} a product that
exists independently in the catalogue.

The four tables were never written as SQL. They were generated by
\texttt{Base.metadata.create\_all(bind=engine)}, which inspects the model
classes at start-up and issues the \texttt{CREATE TABLE} statements.

\begin{figure}[H]
\centering
\shot{0.62\textwidth}{4cm}{fig01-show-tables}{%
Screenshot of the MySQL console after \texttt{mysql -u ormuser -p orm\_workshop}
running \texttt{SHOW TABLES;}. Must show the four tables: \texttt{customers},
\texttt{order\_items}, \texttt{orders}, \texttt{products}.}
\caption{The four tables created from the ORM classes.}
\label{fig:show-tables}
\end{figure}

\section{CRUD Operations}

\begin{table}[H]
\centering\footnotesize
\begin{tabular}{@{}lllll@{}}
\toprule
\textbf{Operation} & \textbf{Method} & \textbf{Customer endpoint} &
\textbf{Product endpoint} & \textbf{ORM operation} \\
\midrule
Create & POST & \texttt{/customers} & \texttt{/products} & \texttt{db.add()} \\
List & GET & \texttt{/customers} & \texttt{/products} & \texttt{select()} \\
Find & GET & \texttt{/customers/\{id\}} & \texttt{/products/\{id\}} & \texttt{db.get()} \\
Update & PUT & \texttt{/customers/\{id\}} & \texttt{/products/\{id\}} & modify object + \texttt{commit()} \\
Delete & DELETE & \texttt{/customers/\{id\}} & \texttt{/products/\{id\}} & \texttt{db.delete()} \\
\bottomrule
\end{tabular}
\caption{CRUD implemented identically for \texttt{Customer} and
\texttt{Product}.}
\end{table}

None of the four operations contains SQL:

\begin{lstlisting}[style=py]
customer = models.Customer(name=data.name, email=data.email)   # create
db.add(customer); db.commit(); db.refresh(customer)

result = db.execute(select(models.Customer))                   # list
customer = db.get(models.Customer, customer_id)                # find by pk

customer.name = data.name                                      # update
db.commit()

db.delete(customer); db.commit()                               # delete
\end{lstlisting}

Three customers were created (Alice Brown, John Smith, Maria Lopez); customer 2
was renamed to ``John Anderson'' and customer 3 was deleted, leaving customers 1
and 2. Four products were created --- Mechanical Keyboard (85.50, stock 20),
Wireless Mouse (28.50, 40), USB-C Hub (42.00, 15) and Webcam (65.00, 10) ---
and the stock of product 3 was updated from 15 to 25.

\begin{figure}[H]
\centering
\begin{minipage}{0.49\textwidth}\centering
\shot{\textwidth}{4.6cm}{fig02-crud-curl}{%
Terminal with the \texttt{curl} commands: the \texttt{POST /customers} creating
Alice Brown with its JSON answer and generated \texttt{id}, and the
\texttt{PUT /customers/2} renaming John with its answer.}
\end{minipage}\hfill
\begin{minipage}{0.49\textwidth}\centering
\shot{\textwidth}{4.6cm}{fig03-crud-mysql}{%
MySQL console running \texttt{SELECT * FROM customers;} and
\texttt{SELECT * FROM products;}: customers 1 and 2 (John already renamed,
Maria absent) and the four products with the USB-C Hub stock at 25.}
\end{minipage}
\caption{Left: creation and update through the API. Right: the same
information stored in MySQL.}
\label{fig:crud}
\end{figure}

\subsection*{Why the program never executes \texttt{INSERT}, \texttt{UPDATE} or
\texttt{DELETE}}

The session acts as a \emph{unit of work}. When an object is loaded or added,
the session keeps a reference to it and a record of its original state; every
later assignment is detected by comparison against that state (dirty tracking).
Nothing is sent at the moment of the assignment. On \texttt{db.commit()} the
session flushes: it collects every new, modified and deleted object, sorts the
operations so that foreign-key dependencies are respected, translates them into
SQL and sends them inside one transaction. That is why
\texttt{customer.name = data.name} followed by \texttt{db.commit()} is enough:
the session noticed that only \texttt{name} changed and produced an
\texttt{UPDATE} affecting only that column. With \texttt{echo=True} the
translation stops being invisible --- every generated statement is printed in
the Uvicorn terminal with its parameters (see Figure~\ref{fig:update-sql}). The
SQL exists; it is simply not written by the programmer.

\section{Relationships and Nested Insert}

A single \texttt{POST /orders} creates one \texttt{Order} and all of its
\texttt{OrderItem} rows. The body sent was
\texttt{\{"customer\_id": 1, "items": [\{"product\_id": 1, "quantity": 2\},
\{"product\_id": 2, "quantity": 1\}, \{"product\_id": 3, "quantity": 3\}]\}}.
The endpoint never assigns a foreign key; it builds objects and links them:

\begin{lstlisting}[style=py]
order = models.Order(customer=customer, status="pending")   # sets orders.customer_id

for item_data in data.items:
    product = db.get(models.Product, item_data.product_id)
    item = models.OrderItem(product=product,                # sets order_items.product_id
                            quantity=item_data.quantity,
                            unit_price=product.price)
    order.items.append(item)                                # sets order_items.order_id

db.add(order)                                               # cascades to the items
db.commit()
\end{lstlisting}

\begin{itemize}
\item Passing \texttt{customer=customer} and \texttt{product=product} associates
\emph{objects}; the ORM records that the corresponding foreign-key columns must
later receive their \texttt{id} values.
\item \texttt{order.items.append(item)} also sets \texttt{item.order}, because
the relationship declares \texttt{back\_populates}, and therefore determines
\texttt{order\_items.order\_id}.
\item \texttt{db.add(order)} adds only the order, but the cascade propagates the
operation to every item.
\item At \texttt{db.commit()} the order still has no \texttt{id}: it is an
\texttt{AUTO\_INCREMENT} generated by MySQL. SQLAlchemy therefore sorts the
inserts by dependency --- first the row in \texttt{orders}, then it reads back
the generated \texttt{id}, and only then the rows in \texttt{order\_items} using
that value. The programmer never handles that identifier.
\end{itemize}

\begin{figure}[H]
\centering
\shot{0.66\textwidth}{5cm}{fig04-nested-insert-mysql}{%
MySQL console running \texttt{SELECT * FROM orders;} and
\texttt{SELECT * FROM order\_items;} right after creating the order. Must show
one row in \texttt{orders} (id 1, \texttt{customer\_id} 1) and three rows in
\texttt{order\_items}, all with \texttt{order\_id} = 1 and \texttt{product\_id}
1, 2 and 3 --- the foreign keys the program never assigned by hand.}
\caption{One order and its three items, created by a single request.}
\label{fig:nested-insert}
\end{figure}

\section{Nested Queries}

The two nested endpoints, \texttt{GET /orders/\{order\_id\}} and
\texttt{GET /customers/\{customer\_id\}/orders}, are implemented in two lines
--- \texttt{order = db.get(models.Order, order\_id)} followed by
\texttt{return order}. A single object is returned; the nesting is produced by
the response schema walking the relationships. This is the real answer of
\texttt{GET /customers/1/orders}, abbreviated to the first item:

\noindent
\begin{minipage}[t]{0.54\textwidth}
\begin{lstlisting}[style=plain]
{
    "id": 1, "name": "Alice Brown",
    "email": "alice@example.com",
    "orders": [
        { "id": 1, "order_date": "2026-09-19",
          "status": "pending",
          "items": [
            { "id": 1, "quantity": 2,
              "unit_price": "85.50",
              "product": {
                  "id": 1,
                  "name": "Mechanical Keyboard",
                  "price": "85.50", "stock": 18 }
            } ]
        } ]
}
\end{lstlisting}
\end{minipage}\hfill
\begin{minipage}[t]{0.43\textwidth}
\vspace{0.2cm}
\begin{lstlisting}[style=plain]
Customer         <- CustomerWithOrders
+-- Order        <- CustomerOrder
    +-- OrderItem<- OrderItemResponse
        +-- Product  <- ProductResponse
\end{lstlisting}
\footnotesize
Four levels out of four separate tables, from one endpoint call.
\normalsize
\end{minipage}

This works because of two mechanisms. The \texttt{relationship()} declarations
turn what in SQL would be a join into attribute access, and those relationships
are loaded lazily: the moment the schema reads \texttt{customer.orders} the
session issues the necessary \texttt{SELECT}, and reading \texttt{item.product}
triggers another. The Pydantic schemas declare the nested shape with
\texttt{from\_attributes=True}, which lets them read from ORM objects. This is
where the split between the two kinds of class pays off: the tables stay flat
and normalised while the API exposes a nested document, and neither has to give
in to the other.

\section{Update and Cascade Delete}

The status of order 1 was changed from \texttt{pending} to \texttt{processing}
with \texttt{order.status = data.status} followed by \texttt{db.commit()}. The
statement observed in the SQLAlchemy output was:

\begin{lstlisting}[style=sql]
UPDATE orders SET status=%(status)s WHERE orders.id = %(orders_id)s
[generated in 0.00099s] {'status': 'processing', 'orders_id': 1}  -->  COMMIT
\end{lstlisting}

Three decisions in that statement were taken by the ORM, not by the program.
\textbf{Which row:} the object came from \texttt{db.get()}, so the session knows
its identity and writes the \texttt{WHERE orders.id = ...} clause itself.
\textbf{Which column:} only \texttt{status} appears in \texttt{SET}, because no
change was detected in the others. \textbf{When:} the assignment alone sends
nothing; the statement is emitted during the flush triggered by
\texttt{db.commit()}. The parameters travel separately from the statement
(\texttt{\%(status)s} is a placeholder), which is what prevents SQL injection.

\begin{figure}[H]
\centering
\shot{0.78\textwidth}{4.2cm}{fig05-update-sql}{%
Uvicorn terminal at the moment of the status update. Must show the
\texttt{UPDATE orders SET status=\%(status)s WHERE orders.id =
\%(orders\_id)s} line, the parameter line
\texttt{\{'status': 'processing', 'orders\_id': 1\}} and the \texttt{COMMIT}.}
\caption{\texttt{UPDATE} generated by a single attribute assignment.}
\label{fig:update-sql}
\end{figure}

Order 1, which had three items, was then deleted with \texttt{db.delete(order)}
and \texttt{db.commit()}. The outcome was: \textbf{(1)} the \texttt{Order} was
deleted and \texttt{orders} was left empty; \textbf{(2)} its three
\texttt{OrderItem} objects were deleted as well, without the program mentioning
them, leaving \texttt{order\_items} empty; \textbf{(3)} the four
\texttt{Product} objects remained untouched, with their stock values unchanged
(18, 39, 22 and 10).

The items disappeared and the products survived because only
\texttt{Order.items} declares \texttt{cascade="all, delete-orphan"}: the
deletion propagates from the order to its items, and SQLAlchemy issues the
\texttt{DELETE} on \texttt{order\_items} first and only then the one on
\texttt{orders}, which is also the order the foreign key requires.
\texttt{Product.order\_items} declares no cascade, so nothing propagates towards
the products.

The underlying reason is semantic. An order item has no independent existence:
``two units of a keyboard'' means nothing outside the order containing it. A
product does: the keyboard stays in the catalogue, available to be sold again,
regardless of how many past orders referenced it. Cascading to products would
destroy the catalogue --- exactly the accident its absence prevents.

\begin{figure}[H]
\centering
\begin{minipage}{0.49\textwidth}\centering
\shot{\textwidth}{4.6cm}{fig06-cascade-before}{%
\textbf{Before} the deletion: \texttt{SELECT * FROM orders;},
\texttt{SELECT * FROM order\_items;} and
\texttt{SELECT id, name, stock FROM products;}. Must show order 1, its three
items and the four products.}
\end{minipage}\hfill
\begin{minipage}{0.49\textwidth}\centering
\shot{\textwidth}{4.6cm}{fig07-cascade-after}{%
\textbf{After} \texttt{DELETE /orders/1}, the same three queries:
\texttt{orders} and \texttt{order\_items} both \texttt{Empty set}, while the
four products are still there unchanged.}
\end{minipage}
\caption{Cascade delete: the items vanished, the products remained.}
\label{fig:cascade}
\end{figure}

\section{Transactions and Rollback}

The \texttt{POST /orders} endpoint was extended so that creating an order also
updates the inventory:

\begin{lstlisting}[style=py]
try:
    for item_data in data.items:
        product = db.get(models.Product, item_data.product_id)
        if product.stock < item_data.quantity:
            raise HTTPException(400, detail=f"Insufficient stock for {product.name}")
        product.stock -= item_data.quantity
        order.items.append(models.OrderItem(product=product,
                                            quantity=item_data.quantity,
                                            unit_price=product.price))
    db.add(order)
    db.commit()
except Exception:
    db.rollback()
    raise
\end{lstlisting}

\texttt{product.stock -= item\_data.quantity} modifies an object loaded by the
session, so it is registered as a pending change exactly like the creation of
the order and of its items. That is the key point: the three kinds of change
accumulate in the same session and therefore belong to the same transaction.
\texttt{db.commit()} flushes every pending change as SQL and confirms the
transaction --- up to that moment nothing is durable, and it either succeeds
completely or fails completely. \texttt{db.rollback()} discards the
transaction, so a half-finished operation leaves no trace.

\subsection*{Successful transaction}

An order was created for customer 2 with two units of product 1 and three of
product 2. The stock of the purchased products decreased by exactly the
requested quantities, in the same transaction that created the order; the other
products were untouched.

\subsection*{Failed transaction}

An order was then requested with one unit of product 1 --- available --- and one
thousand units of product 4, which was not. The answer was HTTP 400 with
\texttt{\{"detail": "Insufficient stock for Webcam"\}}.

\begin{table}[H]
\centering\footnotesize
\begin{tabular}{@{}lccc@{\hspace{1.4cm}}ccc@{}}
\toprule
& \multicolumn{3}{c}{\textbf{Successful transaction}} &
\multicolumn{3}{c}{\textbf{Failed transaction}} \\
\cmidrule(r){2-4}\cmidrule(l){5-7}
\textbf{Product} & Before & Bought & After & Before & Requested & After \\
\midrule
1 --- Mechanical Keyboard & 18 & 2 & \textbf{16} & 16 & 1 (processed) & \textbf{16} \\
2 --- Wireless Mouse & 39 & 3 & \textbf{36} & 36 & --- & 36 \\
3 --- USB-C Hub & 22 & --- & 22 & 22 & --- & 22 \\
4 --- Webcam & 10 & --- & 10 & 10 & 1000 (rejected) & 10 \\
\bottomrule
\end{tabular}
\caption{Inventory before and after each transaction.}
\end{table}

The decisive observation is product 1. The loop processes items in order, so on
reaching the second item it had \emph{already} executed
\texttt{product.stock -= 1} for the keyboard: in memory its stock had gone from
16 to 15. The database nevertheless still reported 16, and no new row appeared
in \texttt{orders}.

This is \textbf{atomicity}: a transaction is an indivisible unit that either
happens entirely or not at all, with no valid intermediate state in which the
customer is charged for a keyboard they never receive. Two things protected the
data. The subtraction only modified an object in memory and marked it pending
--- nothing had been written, because \texttt{db.commit()} was never reached, as
the exception interrupted the loop first. And the \texttt{except} block called
\texttt{db.rollback()}, discarding the transaction along with every pending
change: the subtraction, the order and the item already appended to it.

This is also why the whole loop must live in a single transaction. If each item
were committed separately, the first commit would have made the subtraction on
product 1 durable, and the later failure would have left the inventory
permanently wrong, with one keyboard missing that nobody ever bought.

\begin{figure}[H]
\centering
\begin{minipage}{0.49\textwidth}\centering
\shot{\textwidth}{4.6cm}{fig08-transaction-stock}{%
Terminal showing, in order: the inventory \textbf{before}
(\texttt{curl -s .../products | python -m json.tool}), the \texttt{POST /orders}
for customer 2, and the inventory \textbf{after}. Stock of products 1 and 2
visibly lower (18 to 16, 39 to 36).}
\end{minipage}\hfill
\begin{minipage}{0.49\textwidth}\centering
\shot{\textwidth}{4.6cm}{fig09-rollback}{%
The rollback experiment: inventory before, the \texttt{POST /orders} failing
with \texttt{\{"detail": "Insufficient stock for Webcam"\}} (ideally with the
HTTP 400 visible), and the inventory after. The stock of product 1 must be
\textbf{the same} in both listings.}
\end{minipage}
\caption{Left: successful transaction. Right: failed transaction, with the
partial change on product 1 rolled back.}
\label{fig:transactions}
\end{figure}

\section{Independent Python Client and JSON}

The purpose of the separate client is to show that the API is a real boundary
and not an internal detail. The client imports neither SQLAlchemy nor the
models; it does not know that MySQL exists, nor which tables there are. It could
be rewritten in another language, or run on another machine, without a single
line of the server changing. It communicates through the four HTTP methods:

\begin{lstlisting}[style=py]
requests.get(f"{BASE_URL}/customers")                                   # read
requests.post(f"{BASE_URL}/customers", json=data)                       # create
requests.put(f"{BASE_URL}/orders/{order_id}/status", json=data)         # update
requests.delete(f"{BASE_URL}/orders/{order_id}")                        # delete
\end{lstlisting}

The \texttt{json=} argument serialises the Python dictionary to JSON and sets
the \texttt{Content-Type} header; \texttt{response.json()} parses the answer
back. The file \texttt{client/data/order.json} contains a complete nested order
(customer 1 with three items), which the client reads and forwards:

\begin{lstlisting}[style=py]
with open(filename, "r", encoding="utf-8") as file:
    data = json.load(file)
response = requests.post(f"{BASE_URL}/orders", json=data)
\end{lstlisting}

\texttt{json.load()} returns a Python dictionary with the nested list of items
already converted, and \texttt{requests.post()} serialises it back into the
request body, so the file travels as the payload practically unchanged. On
arrival FastAPI validates it against \texttt{OrderCreate}, which declares
\texttt{items: list[OrderItemCreate]}, and the endpoint creates the order with
its three items in one transaction. A file on disk thus becomes four rows across
two tables, without the client knowing anything about that.

\begin{figure}[H]
\centering
\shot{0.6\textwidth}{5cm}{fig10-client-menu}{%
The Python client running (\texttt{python client/client.py}): the complete menu
and the successful execution of at least one option --- for instance option 5,
``Create order from JSON'', with its \texttt{HTTP: 200} line and the resulting
order's JSON.}
\caption{Independent Python client communicating with the API.}
\label{fig:client}
\end{figure}

\section{Client/Server Test}

The architecture allows the client and the API to run on different computers
(\texttt{Python Client} on A, \texttt{FastAPI + ORM + MySQL} on B), and two
changes are needed, one on each side.

On the \textbf{client}, \texttt{BASE\_URL} must stop pointing at
\texttt{http://127.0.0.1:8000}. The address \texttt{127.0.0.1} is the loopback
interface and always means ``this same machine'', so a client on computer A
using it would look for an API on computer A, where there is none. It must be
replaced by the address of computer B on the local network, e.g.
\texttt{http://192.168.1.50:8000}.

On the \textbf{server}, Uvicorn must be started with \texttt{--host 0.0.0.0}.
That is not an address anyone connects to: it means ``accept connections on
every network interface of this machine''. With the default, Uvicorn binds only
to loopback and any connection arriving from outside is refused before reaching
the application, however correct the client's \texttt{BASE\_URL} may be. MySQL
needs no change: it is only reached by the ORM, which runs on the same machine
as the API, so the only connection crossing the network is the HTTP one.

On this installation the server side was verified --- the process was confirmed
to be listening on \texttt{0.0.0.0:8000} and the API answered through the
machine's network address rather than loopback. The test with a second physical
computer was not carried out, as it requires another machine on the same
network.

\section{Final Exercise}

The endpoint \texttt{GET /products/\{product\_id\}/customers} returns every
customer who purchased a given product (the client got a matching menu option
calling the same URL):

\begin{lstlisting}[style=py]
@app.get("/products/{product_id}/customers",
         response_model=list[schemas.CustomerResponse])
def get_product_customers(product_id: int, db: Session = Depends(get_db)):

    product = db.get(models.Product, product_id)
    if product is None:
        raise HTTPException(status_code=404, detail="Product not found")

    statement = (select(models.Customer)
                 .join(models.Customer.orders)
                 .join(models.Order.items)
                 .where(models.OrderItem.product_id == product_id)
                 .distinct())

    return db.execute(statement).scalars().all()
\end{lstlisting}

\noindent
\begin{minipage}[t]{0.46\textwidth}
\footnotesize\texttt{curl -s http://127.0.0.1:8000/products/1/customers\\
| python -m json.tool}\normalsize
\begin{lstlisting}[style=plain]
[
    { "id": 1, "name": "Alice Brown",
      "email": "alice@example.com" },
    { "id": 2, "name": "John Anderson",
      "email": "john@example.com" }
]
\end{lstlisting}
\end{minipage}\hfill
\begin{minipage}[t]{0.51\textwidth}
\vspace{0.1cm}
\footnotesize
Both customers bought the Mechanical Keyboard: John in order 2 and Alice in
order 3. As a contrast, \texttt{/products/4/customers} returns an empty list,
because the Webcam was never sold successfully --- the only order that requested
it is precisely the one that was rolled back.
\normalsize
\end{minipage}

\medskip
Customers and products are not directly connected: two tables lie between them
(\texttt{Customer --< Order --< OrderItem >-- Product}).
\texttt{.join(models.Customer.orders)} takes the first step, from
\texttt{customers} to \texttt{orders}; \texttt{.join(models.Order.items)} takes
the second, to \texttt{order\_items}. Both joins are expressed with a
\emph{relationship} rather than an ON condition: the ORM already knows the
linking columns and writes them itself. The \texttt{where} clause then keeps
only the lines referring to the product. The fourth table is never joined,
because \texttt{order\_items.product\_id} already holds the value being
compared, which saves one join.

\begin{lstlisting}[style=sql]
SELECT DISTINCT customers.id, customers.name, customers.email
FROM customers
INNER JOIN orders ON customers.id = orders.customer_id
INNER JOIN order_items ON orders.id = order_items.order_id
WHERE order_items.product_id = %(product_id_1)s
\end{lstlisting}

\texttt{.distinct()} is necessary because a join multiplies rows: the result has
one row \emph{per order item}, not per customer, with the customer's columns
repeated in each. A customer who bought the same product in three orders would
appear three times, and the same happens if one order contains two lines for the
same product. \texttt{DISTINCT} collapses those duplicates into one entry per
customer. The duplication is not a bug in the query --- it is the natural
behaviour of a join, and \texttt{DISTINCT} is the standard way of adapting that
result to a question asked about customers rather than about purchase lines.

\begin{figure}[H]
\centering
\shot{0.6\textwidth}{4.2cm}{fig11-final-exercise}{%
Terminal running \texttt{curl -s http://127.0.0.1:8000/products/1/customers |
python -m json.tool} with its answer, the two customers. If it fits, include the
call to \texttt{/products/4/customers} returning \texttt{[]}, which shows the
endpoint does not simply list every customer.}
\caption{Customers who purchased product 1.}
\label{fig:final}
\end{figure}

\section{Discussion Questions}

\begin{enumerate}[leftmargin=0.7cm]

\item \textbf{What is an ORM, and what advantages does it provide over writing
SQL directly?} A layer that maps classes to tables and instances to rows, so the
program works with objects instead of result sets. The schema is defined once in
the models and the tables are generated from it; relationships turn joins into
attribute access, which is how a four-level nested response came out of a single
\texttt{db.get()}; every query is parametrised by construction, removing a whole
family of injection bugs; and the session groups changes into transactions,
which made the rollback work almost for free. The price is a loss of control:
the generated SQL is not always the one an expert would write, and problems such
as the N+1 query pattern appear only when the \texttt{echo=True} output is
actually read.

\item \textbf{Difference between a SQLAlchemy model and a Pydantic schema?} They
answer different questions. A model describes \emph{how data is stored}: tables,
columns, SQL types, foreign keys. A schema describes \emph{how data travels}:
the shape of the JSON the API accepts and returns. Separating them is what lets
the tables stay flat while the API exposes a nested document, and lets the two
differ on purpose --- \texttt{CustomerCreate} requires name and email,
\texttt{CustomerUpdate} makes both optional, \texttt{CustomerResponse} adds the
\texttt{id} that only exists after the insert.

\item \textbf{Purpose of \texttt{db.commit()}?} To confirm the transaction. It
first flushes --- translating every pending change into SQL and sending it in an
order that respects foreign-key dependencies --- and then commits, which is the
point at which the changes become durable and visible to other connections.
Before that call, everything exists only in the session.

\item \textbf{Why is \texttt{db.rollback()} important?} Because it guarantees
that a failed operation leaves no trace. Without it the partial changes of an
interrupted operation --- here, the stock already subtracted before the failure
was detected --- would remain in the session and could be written by a later
commit. It discards the transaction and returns the session to a clean state,
which is what makes a failure safe.

\item \textbf{How does SQLAlchemy represent a one-to-many relationship?} With
two complementary pieces: a \texttt{ForeignKey} column on the many side
(\texttt{orders.customer\_id}), which is the actual link in the database, and a
pair of \texttt{relationship()} declarations joined by \texttt{back\_populates}
(\texttt{Customer.orders} as a list, \texttt{Order.customer} as a single
object), which give navigation in both directions and stay consistent when
either is modified.

\item \textbf{What does \texttt{cascade="all, delete-orphan"} accomplish?} It
propagates session operations from a parent to its children. \texttt{all} makes
operations such as \texttt{add} and \texttt{delete} reach the children --- which
is why \texttt{db.add(order)} also saved its items and deleting the order also
deleted them. \texttt{delete-orphan} adds a second rule: a child removed from
the parent's collection is deleted, because it cannot exist on its own.

\item \textbf{Why does deleting an Order remove its OrderItems but not its
Products?} Because only \texttt{Order.items} declares a cascade, and that
reflects ownership. An order owns its items: a line has no meaning outside the
order containing it. An order only references products: the product exists
independently, before the order and after it. Declaring the cascade on the
product side would destroy part of the catalogue on every deletion.

\item \textbf{What advantage does a nested insert provide?} It makes the whole
operation a single unit in three senses: one request instead of one per item for
the client; one transaction for the database, so an order can never be left
without its items; and no intermediate identifiers for the programmer, since the
order's \texttt{id} does not exist until the insert happens and SQLAlchemy reads
it back and propagates it.

\item \textbf{Why should an Order and its inventory changes be one
transaction?} Because they describe a single business fact: this sale happened.
Splitting them allows states no rule of the domain admits --- an order without
its stock reserved, or stock subtracted for an order never recorded. In one
transaction the alternatives are exactly two: either the order exists and the
inventory reflects it, or nothing happened.

\item \textbf{Role of FastAPI and role of SQLAlchemy?} FastAPI is the boundary
facing the outside world: it receives HTTP requests, routes them, validates the
JSON against the schemas, serialises the response and handles errors ---
\texttt{HTTPException(400)} is what turned an insufficient-stock condition into
a well-formed HTTP answer. SQLAlchemy is the boundary facing the database: it
maps classes to tables, tracks changes, generates the SQL and manages the
transaction. They never overlap: FastAPI knows nothing about SQL and SQLAlchemy
nothing about HTTP. The endpoint functions are the thin layer where they meet.

\end{enumerate}

\section{Conclusions}

\begin{enumerate}[leftmargin=0.7cm]

\item \textbf{The ORM does not hide SQL, it generates it --- and the distinction
matters.} Reading the \texttt{echo=True} output showed that behind
\texttt{customer.name = "John Anderson"} there is an \texttt{UPDATE} touching
exactly one column of one row, and that behind a nested insert there is an
ordered sequence of statements with the generated identifier read back in
between. The productive way to use an ORM is not to forget SQL but to know which
SQL each object operation produces; otherwise it is impossible to notice, for
instance, that the nested query triggers one \texttt{SELECT} per relationship
traversed.

\item \textbf{Relationships are where the object model stops being a copy of the
tables.} The tables are flat and normalised, yet the response of
\texttt{/customers/1/orders} has four levels of nesting, and nothing in the
database changed to make that possible: it came from \texttt{relationship()}
turning joins into attribute access and from the schemas describing the desired
shape. Storage model and exposed model are two different designs, and separating
them is what keeps each free to change.

\item \textbf{Cascade is a modelling decision, not a configuration option.}
Declaring \texttt{cascade="all, delete-orphan"} on \texttt{Order.items} and
omitting it on \texttt{Product.order\_items} encodes which objects own others
and which merely reference them. The delete experiment made the consequence
visible: three items disappeared without being named while four products
survived. Getting that decision wrong produces no error message --- it silently
destroys data.

\item \textbf{Atomicity is only visible when something fails.} The successful
transaction proves little; committing each item separately would look identical.
The failed one is where the design shows itself: the stock of product 1 had
already been subtracted in memory when the exception was raised for product 4,
and it still ended up unchanged in the database, because the subtraction never
crossed a \texttt{commit()} and \texttt{rollback()} discarded it. What is
protected is not the individual operation but the invariant connecting them: an
order and its inventory effect are one fact.

\item \textbf{The separation between API, ORM and client makes the system
replaceable in parts.} The client never imports a model, never opens a session
and does not know MySQL exists. Because of that, the same API answers through
loopback and through the machine's network address, and moving the client to
another computer requires changing exactly one constant. Each layer can be
replaced without touching the others, and that is a direct consequence of none
of them knowing the internals of its neighbours.

\end{enumerate}

\end{document}
